% Generated by tools/make_body.py; edit generator, not this file.
% UNPROOFREAD OCR TRANSCRIPTION: page-image verification remains required.
\chapter{Complex Functions}\label{chap:2}

% Original PDF physical page 37
\sourcepage{37}
\setcounter{section}{0}\section{Introduction To The Concept Of Analytic Function}\label{sec:2.1}
The theory of functions of a complex variable aims at extending calculus to the complex domain. Both differentiation and integration acquire new depth and significance; at the same time the range of applicability becomes radically restricted. Indeed, only the analytic or holomorphic functions\index{Function!holomorphic}\index{Holomorphic function} can be freely differentiated and integrated. They are the only true “functions” in the sense of the French “Théorie des fonctions” or the German “Funktionentheorie.”

Nevertheless, we shall use the term “function” in its modern meaning. Therefore, when stepping up to complex numbers we have to consider four different kinds of functions: real functions of a real variable, real functions of a complex variable, complex functions\index{Complex function|(}\index{Function!complex|(}\index{Complex function}\index{Function!complex} of a real variable, and complex functions of a complex variable. As a practical matter we agree that the letters z and w shall always denote complex variables; thus, to indicate a complex function of a complex variable we use the notation $w = f(z)$.\hyperlink{fn:37:dagger}{\textsuperscript{\ensuremath{\dagger}}} The notation $y = f(x)$ will be used in a neutral manner with the understanding that x and y can be either real or complex. When we want to indicate that a variable is definitely restricted to real values, we shall usually denote it by t. By these agreements we
\begingroup\renewcommand{\thefootnote}{}\footnotetext{\hypertarget{fn:37:dagger}{}\textsuperscript{\ensuremath{\dagger}} Modern students are well aware that $f$ stands for the function and $f(z)$ for a value of the function. However, analysts are traditionally minded and continue to speak of “the function $f(z)$.”}\endgroup
% Original PDF physical page 38
\sourcepage{38}
do not wish to cancel the earlier convention whereby a notation $z = x + iy$ automatically implies that x and y are real.

It is essential that the law by which a function is defined be formulated in clear and unambiguous terms. In other words, all functions must be well defined and consequently, until further notice, single-valued.\hyperlink{fn:38:dagger}{\textsuperscript{\ensuremath{\dagger}}}

It is not necessary that a function be defined for all values of the independent variable. For the moment we shall deliberately underemphasize the role of point set theory. Therefore we make merely an informal agreement that every function be defined on an open set, by which we mean that if $f(a)$ is defined, then $f(x)$ is defined for all x sufficiently close to a. The formal treatment of point set topology is deferred until the next chapter.

\setcounter{section}{1}\setcounter{subsection}{0}\subsection{Limits and Continuity}\label{sec:2.1.1}
The following basic definition will be adopted:

\begin{originalstatement}{Definition}{1}\label{def:2.1}
The function $f(x)$ is said to have the limit\index{Limit|(}\index{Limit} $A$ as $x$ tends to $a$,
\begin{equation*}
\lim _ {x \rightarrow a} f (x) = A,\tag{1}\label{eq:2.1}
\end{equation*}
if and only if the following is true:

For every $\varepsilon > 0$ there exists a number $\delta > 0$ with the property that $|f(x) - A| < \varepsilon$ for all values of $x$ such that $|x - a| < \delta$ and $x \neq a$.

\end{originalstatement}
This definition makes decisive use of the absolute value. Since the notion of absolute value has a meaning for complex as well as for real numbers, we can use the same definition regardless of whether the variable x and the function $f(x)$ are real or complex.

As an alternative simpler notation we sometimes write: $f(x) \to A$ for $x \to a$.

There are some familiar variants of the definition which correspond to the case where a or A is infinite. In the real case we can distinguish between the limits $+\infty$ and $-\infty$, but in the complex case there is only one infinite limit. We trust the reader to formulate correct definitions to cover all the possibilities.

The well-known results concerning the limit of a sum, a product, and a quotient continue to hold in the complex case. Indeed, the proofs depend only on the properties of the absolute value expressed by
\begin{equation*}
| a b | = | a | \cdot | b | \quad \text { and } \quad | a + b | \leq | a | + | b |.
\end{equation*}
\begingroup\renewcommand{\thefootnote}{}\footnotetext{\hypertarget{fn:38:dagger}{}\textsuperscript{\ensuremath{\dagger}} We shall sometimes use the pleonastic term \emph{single-valued function}\index{Function!single-valued}\index{Single-valued function} to underline that the function has only one value for each value of the variable.}\endgroup

% Original PDF physical page 39
\sourcepage{39}
Condition \hyperref[eq:2.1]{(1)} is evidently equivalent to
\begin{equation*}
\lim _ {x \rightarrow a} \overline {{{{f (x)}}}} = \bar {A}.\tag{2}\label{eq:2.2}
\end{equation*}
From \hyperref[eq:2.1]{(1)} and \hyperref[eq:2.2]{(2)} we obtain
\begin{equation*}
\begin{array}{l} \lim _ {x \to a} \operatorname{Re} f (x) = \operatorname{Re} A \\ \lim _ {x \to a} \operatorname{Im} f (x) = \operatorname{Im} A. \end{array}\tag{3}\label{eq:2.3}
\end{equation*}
Conversely, \hyperref[eq:2.1]{(1)} is a consequence of \hyperref[eq:2.3]{(3)}.

The function $f(x)$ is said to be continuous at $a$ if and only if $\lim_{x \to a} f(x) = f(a)$. A continuous function\index{Continuous function}\index{Function!continuous}, without further qualification, is one which is continuous at all points where it is defined.

The sum $f(x) + g(x)$ and the product $f(x)g(x)$ of two continuous functions are continuous; the quotient $f(x)/g(x)$ is defined and continuous at a if and only if $g(a) \neq 0$. If $f(x)$ is continuous, so are $\operatorname{Re} f(x)$, $\operatorname{Im} f(x)$, and $|f(x)|$.

The derivative\index{Derivative|(}\index{Derivative} of a function is defined as a particular limit\index{Limit} and can be considered regardless of whether the variables are real or complex. The formal definition is
\begin{equation*}
f ^ {\prime} (a) = \lim _ {x \rightarrow a} \frac {f (x) - f (a)}{x - a}.\tag{4}\label{eq:2.4}
\end{equation*}
The usual rules for forming the derivative of a sum, a product, or a quotient are all valid. The derivative of a composite function is determined by the chain rule.

There is nevertheless a fundamental difference between the cases of a real and a complex independent variable. To illustrate our point, let $f(z)$ be a real function of a complex variable whose derivative exists at z = a. Then $f'(a)$ is on one side real, for it is the limit of the quotients
\begin{equation*}
\frac {f (a + h) - f (a)}{h}
\end{equation*}
as $h$ tends to zero through real values. On the other side it is also the limit of the quotients
\begin{equation*}
\frac {f (a + i h) - f (a)}{i h}
\end{equation*}
and as such purely imaginary. Therefore $f'(a)$ must be zero. Thus a real function of a complex variable either has the derivative zero, or else the derivative does not exist.


% Original PDF physical page 40
\sourcepage{40}
The case of a complex function\index{Complex function}\index{Function!complex} of a real variable can be reduced to the real case. If we write $z(t) = x(t) + iy(t)$ we find indeed
\begin{equation*}
z ^ {\prime} (t) = x ^ {\prime} (t) + i y ^ {\prime} (t),
\end{equation*}
and the existence of $z'(t)$ is equivalent to the simultaneous existence of $x'(t)$ and $y'(t)$. The complex notation has nevertheless certain formal advantages which it would be unwise to give up.

In contrast, the existence of the derivative\index{Derivative} of a complex function of a complex variable has far-reaching consequences for the structural properties of the function. The investigation of these consequences is the central theme in complex-function theory.

\setcounter{section}{1}\setcounter{subsection}{1}\subsection{Analytic Functions}\label{sec:2.1.2}
The class of analytic functions\index{Function!analytic|(}\index{Function!analytic} is formed by the complex functions of a complex variable which possess a derivative wherever the function is defined. The term holomorphic function\index{Function!holomorphic}\index{Holomorphic function} is used with identical meaning. For the purpose of this preliminary investigation the reader may think primarily of functions which are defined in the whole plane.

The sum and the product of two analytic functions are again analytic. The same is true of the quotient $f(z)/g(z)$ of two analytic functions, provided that $g(z)$ does not vanish. In the general case it is necessary to exclude the points at which $g(z)=0$. Strictly speaking, this very typical case will thus not be included in our considerations, but it will be clear that the results remain valid except for obvious modifications.

The definition of the derivative can be rewritten in the form
\begin{equation*}
f ^ {\prime} (z) = \lim _ {h \rightarrow 0} \frac {f (z + h) - f (z)}{h}.
\end{equation*}
As a first consequence $f(z)$ is necessarily continuous. Indeed, from $f(z + h) - f(z) = h \cdot (f(z + h) - f(z)) / h$ we obtain
\begin{equation*}
\lim _ {h \rightarrow 0} (f (z + h) - f (z)) = 0 \cdot f ^ {\prime} (z) = 0.
\end{equation*}
If we write $f(z) = u(z) + iv(z)$ it follows, moreover, that $u(z)$ and $v(z)$ are both continuous.

The limit\index{Limit} of the difference quotient must be the same regardless of the way in which h approaches zero. If we choose real values for h, then the imaginary part y is kept constant, and the derivative becomes a partial derivative with respect to x. We have thus
\begin{equation*}
f ^ {\prime} (z) = \frac {\partial f}{\partial x} = \frac {\partial u}{\partial x} + i \frac {\partial v}{\partial x}.
\end{equation*}
% Original PDF physical page 41
\index{Derivative|)}\index{Limit|)}\sourcepage{41}
Similarly,\index{Conjugate harmonic function|(}\index{Function!conjugate harmonic|(} if we substitute purely imaginary values $ik$ for $h$, we obtain
\begin{equation*}
f ^ {\prime} (z) = \lim _ {k \rightarrow 0} \frac {f (z + i k) - f (z)}{i k} = - i \frac {\partial f}{\partial y} = - i \frac {\partial u}{\partial y} + \frac {\partial v}{\partial y}.
\end{equation*}
It follows that $f(z)$ must satisfy the partial differential equation
\begin{equation*}
\frac {\partial f}{\partial x} = - i \frac {\partial f}{\partial y}\tag{5}\label{eq:2.5}
\end{equation*}
which resolves into the real equations
\begin{equation*}
\frac {\partial u}{\partial x} = \frac {\partial v}{\partial y}, \quad \frac {\partial u}{\partial y} = - \frac {\partial v}{\partial x}.\tag{6}\label{eq:2.6}
\end{equation*}
These are the Cauchy-Riemann differential equations which must be satisfied by the real and imaginary part of any analytic function\index{Function!analytic}.\hyperlink{fn:41:dagger}{\textsuperscript{\ensuremath{\dagger}}}

We remark that the existence of the four partial derivatives in \hyperref[eq:2.6]{(6)} is implied by the existence of $f'(z)$. Using \hyperref[eq:2.6]{(6)} we can write down four formally different expressions for $f'(z)$; the simplest is
\begin{equation*}
f ^ {\prime} (z) = \frac {\partial u}{\partial x} + i \frac {\partial v}{\partial x}.
\end{equation*}
For the quantity $|f'(z)|^2$ we have, for instance,
\begin{equation*}
\left| f ^ {\prime} (z) \right| ^ {2} = \left(\frac {\partial u}{\partial x}\right) ^ {2} + \left(\frac {\partial u}{\partial y}\right) ^ {2} = \left(\frac {\partial u}{\partial x}\right) ^ {2} + \left(\frac {\partial v}{\partial x}\right) ^ {2} = \frac {\partial u}{\partial x} \frac {\partial v}{\partial y} - \frac {\partial u}{\partial y} \frac {\partial v}{\partial x}.
\end{equation*}
The last expression shows that $|f'(z)|^2$ is the Jacobian\index{Jacobian} of $u$ and $v$ with respect to $x$ and $y$.

We shall prove later that the derivative of an analytic function is itself analytic. By this fact u and v will have continuous partial derivatives of all orders, and in particular the mixed derivatives will be equal. Using this information we obtain from \hyperref[eq:2.6]{(6)}
\begin{equation*}
\begin{array}{l} \Delta u = \frac {\partial^ {2} u}{\partial x ^ {2}} + \frac {\partial^ {2} u}{\partial y ^ {2}} = 0 \\ \Delta v = \frac {\partial^ {2} v}{\partial x ^ {2}} + \frac {\partial^ {2} v}{\partial y ^ {2}} = 0. \end{array}
\end{equation*}
A function u which satisfies Laplace's equation\index{Laplace's equation} $\Delta u = 0$ is said to be harmonic. The real and imaginary part of an analytic function are thus harmonic. If two harmonic functions\index{Function!harmonic}\index{Harmonic function} u and v satisfy the Cauchy-Riemann equations \hyperref[eq:2.6]{(6)}, then v is said to be the conjugate harmonic func%
\begingroup\renewcommand{\thefootnote}{}\footnotetext{\hypertarget{fn:41:dagger}{}\textsuperscript{\ensuremath{\dagger}} Augustin Cauchy\index{Cauchy A@Cauchy, A.|hyperpagenote} (1789–1857) and Bernhard Riemann\index{Riemann B@Riemann, B.|hyperpagenote} (1826–1866) are regarded as the founders of complex-function theory. Riemann's work emphasized the geometric aspects in contrast to the purely analytic approach of Cauchy.}\endgroup%
% Original PDF physical page 42
\sourcepage{42}%
tion of u. Actually, v is determined only up to an additive constant, so that the use of the definite article, although traditional, is not quite accurate. In the same sense, u is the conjugate harmonic function\index{Conjugate harmonic function}\index{Function!conjugate harmonic} of -v.

This is not the place to discuss the weakest conditions of regularity which can be imposed on harmonic functions. We wish to prove, however, that the function $u + iv$ determined by a pair of conjugate harmonic functions is always analytic, and for this purpose we make the explicit assumption that u and v have continuous first-order partial derivatives. It is proved in calculus, under exactly these regularity conditions, that we can write
\begin{equation*}
u (x + h, y + k) - u (x, y) = \frac {\partial u}{\partial x} h + \frac {\partial u}{\partial y} k + \varepsilon_ {1}
\end{equation*}
\begin{equation*}
v (x + h, y + k) - v (x, y) = \frac {\partial v}{\partial x} h + \frac {\partial v}{\partial y} k + \varepsilon_ {2},
\end{equation*}
where the remainders $\varepsilon_{1}, \varepsilon_{2}$ tend to zero more rapidly than $h + ik$ in the sense that $\varepsilon_{1}/(h + ik) \to 0$ and $\varepsilon_{2}/(h + ik) \to 0$ for $h + ik \to 0$. With the notation $f(z) = u(x,y) + iv(x,y)$ we obtain by virtue of the relations \hyperref[eq:2.6]{(6)}
\begin{equation*}
f (z + h + i k) - f (z) = \left(\frac {\partial u}{\partial x} + i \frac {\partial v}{\partial x}\right) (h + i k) + \varepsilon_ {1} + i \varepsilon_ {2}
\end{equation*}
and hence
\begin{equation*}
\lim _ {h + i k \rightarrow 0} \frac {f (z + h + i k) - f (z)}{h + i k} = \frac {\partial u}{\partial x} + i \frac {\partial v}{\partial x}.
\end{equation*}
We conclude that $f(z)$ is analytic.

If $u(x,y)$ and $v(x,y)$ have continuous first-order partial derivatives which satisfy the Cauchy-Riemann differential equations, then $f(z) = u(z) + iv(z)$ is analytic with continuous derivative $f'(z)$, and conversely.

The conjugate of a harmonic function can be found by integration, and in simple cases the computation can be made explicit. For instance, $u = x^{2} - y^{2}$ is harmonic and $\partial u/\partial x = 2x$, $\partial u/\partial y = -2y$. The conjugate function must therefore satisfy
\begin{equation*}
\frac {\partial v}{\partial x} = 2 y, \quad \frac {\partial v}{\partial y} = 2 x.
\end{equation*}
From the first equation $v = 2xy + \varphi(y)$, where $\varphi(y)$ is a function of y alone. Substitution in the second equation yields $\varphi'(y) = 0$. Hence $\varphi(y)$ is a constant, and the most general conjugate function of $x^{2} - y^{2}$ is $2xy + c$ where c is a constant. Observe that $x^{2} - y^{2} + 2ixy = z^{2}$. The analytic function\index{Function!analytic} with the real part $x^{2} - y^{2}$ is hence $z^{2} + ic$.


% Original PDF physical page 43
\index{Conjugate harmonic function|)}\index{Function!conjugate harmonic|)}\sourcepage{43}
There is an interesting formal procedure which throws considerable light on the nature of analytic functions\index{Function!analytic}. We present this procedure with an explicit warning to the reader that it is purely formal and does not possess any power of proof.

Consider a complex function\index{Complex function}\index{Function!complex} $f(x,y)$ of two real variables. Introducing the complex variable $z = x + iy$ and its conjugate $\bar{z} = x - iy$, we can write $x = \frac{1}{2}(z + \bar{z})$, $y = -\frac{1}{2}i(z - \bar{z})$. With this change of variable we can consider $f(x,y)$ as a function of z and $\bar{z}$ which we will treat as independent variables (forgetting that they are in fact conjugate to each other). If the rules of calculus were applicable, we would obtain
\begin{equation*}
\frac {\partial f}{\partial z} = \frac {1}{2} \left(\frac {\partial f}{\partial x} - i \frac {\partial f}{\partial y}\right), \quad \frac {\partial f}{\partial \bar {z}} = \frac {1}{2} \left(\frac {\partial f}{\partial x} + i \frac {\partial f}{\partial y}\right).
\end{equation*}
These expressions have no convenient definition as limits, but we can nevertheless introduce them as symbolic derivatives\index{Symbolic derivative} with respect to z and $\bar{z}$. By comparison with \hyperref[eq:2.5]{(5)} we find that analytic functions are characterized by the condition $\partial f/\partial\bar{z}=0$. We are thus tempted to say that an analytic function is independent of $\bar{z}$, and a function of z alone.

This formal reasoning supports the point of view that analytic functions are true functions of a complex variable as opposed to functions which are more adequately described as complex functions of two real variables.

By similar formal arguments we can derive a very simple method which allows us to compute, without use of integration, the analytic function $f(z)$ whose real part is a given harmonic function $u(x,y)$. We remark first that the conjugate function $\overline{f(z)}$ has the derivative zero with respect to z and may, therefore, be considered as a function of $\bar{z}$; we denote this function by $\bar{f}(\bar{z})$. With this notation we can write down the identity
\begin{equation*}
u (x, y) = \frac {1}{2} [ f (x + i y) + \bar {f} (x - i y) ].
\end{equation*}
It is reasonable to expect that this is a formal identity, and then it holds even when x and y are complex. If we substitute x = z/2, y = z/2i, we obtain
\begin{equation*}
u (z / 2, z / 2 i) = \frac {1}{2} [ f (z) + \bar {f} (0) ].
\end{equation*}
Since $f(z)$ is only determined up to a purely imaginary constant, we may as well assume that $f(0)$ is real, which implies $\bar{f}(0) = u(0,0)$. The function $f(z)$ can thus be computed by means of the formula
\begin{equation*}
f (z) = 2 u (z / 2, z / 2 i) - u (0, 0).
\end{equation*}
A purely imaginary constant can be added at will.

In this form the method is definitely limited to functions $u(x,y)$ which
% Original PDF physical page 44
\sourcepage{44}
are rational in x and y, for the function must have a meaning for complex values of the argument. Suffice it to say that the method can be extended to the general case and that a complete justification can be given.

\subsection*{EXERCISES}
\originalnumber{1}{exercise:2.1.2.1.1}
1. If $g(w)$ and $f(z)$ are analytic functions\index{Function!analytic}, show that $g(f(z))$ is also analytic.

\originalnumber{2}{exercise:2.1.2.1.2}
2. Verify Cauchy-Riemann's equations for the functions $z^2$ and $z^3$.

\originalnumber{3}{exercise:2.1.2.1.3}
3. Find the most general harmonic polynomial\index{Polynomial|(}\index{Polynomial} of the form $ax^{3} + bx^{2}y + cxy^{2} + dy^{3}$. Determine the conjugate harmonic function and the corresponding analytic function by integration and by the formal method.

\originalnumber{4}{exercise:2.1.2.1.4}
4. Show that an analytic function cannot have a constant absolute value without reducing to a constant.

\originalnumber{5}{exercise:2.1.2.1.5}
5. Prove rigorously that the functions $f(z)$ and $\overline{f(\bar{z})}$ are simultaneously analytic.

\originalnumber{6}{exercise:2.1.2.1.6}
6. Prove that the functions $u(z)$ and $u(\bar{z})$ are simultaneously harmonic.

\originalnumber{7}{exercise:2.1.2.1.7}
7. Show that a harmonic function satisfies the formal differential equation
\begin{equation*}
\frac {\partial^ {2} u}{\partial z \partial \bar {z}} = 0.
\end{equation*}
\setcounter{section}{1}\setcounter{subsection}{2}\subsection{Polynomials}\label{sec:2.1.3}
Every constant is an analytic function with the derivative 0. The simplest nonconstant analytic function is z whose derivative is 1. Since the sum and product of two analytic functions are again analytic, it follows that every polynomial
\begin{equation*}
P (z) = a _ {0} + a _ {1} z + \dots + a _ {n} z ^ {n}\tag{7}\label{eq:2.7}
\end{equation*}
is an analytic function. Its derivative is
\begin{equation*}
P ^ {\prime} (z) = a _ {1} + 2 a _ {2} z + \dots + n a _ {n} z ^ {n - 1}.
\end{equation*}
The notation \hyperref[eq:2.7]{(7)} shall imply that $a_{n} \neq 0$, and the polynomial is then said to be of degree n. The constant 0, considered as a polynomial, is in many respects exceptional and will be excluded from our considerations.\hyperlink{fn:44:dagger}{\textsuperscript{\ensuremath{\dagger}}}

For n > 0 the equation $P(z) = 0$ has at least one root. This is the so-called fundamental theorem of algebra\index{Fundamental theorem of algebra} which we shall prove later. If $P(\alpha_{1}) = 0$, it is shown in elementary algebra that $P(z) = (z - \alpha_{1})P_{1}(z)$ where $P_{1}(z)$ is a polynomial of degree n - 1. Repetition of this process finally leads to a complete factorization
\begin{equation*}
P (z) = a _ {n} \left(z - \alpha_ {1}\right) \left(z - \alpha_ {2}\right) \cdot \cdot \cdot \left(z - \alpha_ {n}\right)\tag{8}\label{eq:2.8}
\end{equation*}
\begingroup\renewcommand{\thefootnote}{}\footnotetext{\hypertarget{fn:44:dagger}{}\textsuperscript{\ensuremath{\dagger}} For formal reasons, if the constant 0 is regarded as a polynomial, its degree is set equal to $-\infty$.}\endgroup

% Original PDF physical page 45
\index{Function!analytic|)}\sourcepage{45}
where\index{Order, algebraic!of zero|(} the $\alpha_{1}, \alpha_{2}, \ldots, \alpha_{n}$ are not necessarily distinct. From the factorization we conclude that $P(z)$ does not vanish for any value of z different from $\alpha_{1}, \alpha_{2}, \ldots, \alpha_{n}$. Moreover, the factorization is uniquely determined except for the order of the factors.

If exactly h of the $\alpha_{j}$ coincide, their common value is called a zero\index{Zero} of $P(z)$ of the order h. We find that the sum of the orders of the zeros of a polynomial\index{Polynomial} is equal to its degree. More simply, if each zero is counted as many times as its order indicates, a polynomial of degree n has exactly n zeros.

The order of a zero $\alpha$ can also be determined by consideration of the successive derivatives of $P(z)$ for $z = \alpha$. Suppose that $\alpha$ is a zero of order h. Then we can write $P(z) = (z - \alpha)^{h}P_{h}(z)$ with $P_{h}(\alpha) \neq 0$. Successive derivation yields $P(\alpha) = P'(\alpha) = \cdots = P^{(h-1)}(\alpha) = 0$ while $P^{(h)}(\alpha) \neq 0$. In other words, the order of a zero equals the order of the first nonvanishing derivative. A zero of order 1 is called a simple zero and is characterized by the conditions $P(\alpha) = 0$, $P'(\alpha) \neq 0$.

As an application we shall prove the following theorem, known as Lucas's theorem\index{Lucas's theorem}:

\begin{originalstatement}{Theorem}{1}\label{thm:2.1}
If all zeros of a polynomial $P(z)$ lie in a half plane, then all zeros of the derivative $P'(z)$ lie in the same half plane.
\end{originalstatement}
From \hyperref[eq:2.8]{(8)} we obtain
\begin{equation*}
\frac {P ^ {\prime} (z)}{P (z)} = \frac {1}{z - \alpha_ {1}} + \dots + \frac {1}{z - \alpha_ {n}}.\tag{9}\label{eq:2.9}
\end{equation*}
Suppose that the half plane H is defined as the part of the plane where $\operatorname{Im}(z-a)/b<0$ (see Chap.~\hyperref[chap:1]{1}, Sec.~\hyperref[sec:1.2.3]{2.3}). If $\alpha_{k}$ is in H and z is not, we have then
\begin{equation*}
\operatorname{Im} \frac {z - \alpha_ {k}}{b} = \operatorname{Im} \frac {z - a}{b} - \operatorname{Im} \frac {\alpha_ {k} - a}{b} > 0.
\end{equation*}
But the imaginary parts of reciprocal numbers have opposite sign. Therefore, under the same assumption, $\operatorname{Im} b(z - \alpha_{k})^{-1} < 0$. If this is true for all k we conclude from \hyperref[eq:2.9]{(9)} that
\begin{equation*}
\operatorname{Im} \frac {b P ^ {\prime} (z)}{P (z)} = \sum_ {k = 1} ^ {n} \operatorname{Im} \frac {b}{z - \alpha_ {k}} <   0,
\end{equation*}
and consequently $P'(z) \neq 0$.

In a sharper formulation the theorem tells us that the smallest convex polygon that contains the zeros of $P(z)$ also contains the zeros of $P'(z)$.


% Original PDF physical page 46
\index{Order, algebraic!of zero|)}\index{Polynomial|)}\sourcepage{46}
\setcounter{section}{1}\setcounter{subsection}{3}\subsection{Rational Functions}\index{Order, algebraic!of a pole|(}\label{sec:2.1.4}
We turn to the case of a rational function\index{Rational function|(}\index{Rational function}
\begin{equation*}
R (z) = \frac {P (z)}{Q (z)},\tag{10}\label{eq:2.10}
\end{equation*}
given as the quotient of two polynomials. We assume, and this is essential, that $P(z)$ and $Q(z)$ have no common factors and hence no common zeros. $R(z)$ will be given the value $\infty$ at the zeros of $Q(z)$. It must therefore be considered as a function with values in the extended plane, and as such it is continuous. The zeros of $Q(z)$ are called poles\index{Pole} of $R(z)$, and the order of a pole is by definition equal to the order of the corresponding zero of $Q(z)$.

The derivative
\begin{equation*}
R ^ {\prime} (z) = \frac {P ^ {\prime} (z) Q (z) - Q ^ {\prime} (z) P (z)}{Q (z) ^ {2}}\tag{11}\label{eq:2.11}
\end{equation*}
exists only when $Q(z) \neq 0$. However, as a rational function defined by the right-hand member of \hyperref[eq:2.11]{(11)}, $R'(z)$ has the same poles as $R(z)$, the order of each pole being increased by one. In case $Q(z)$ has multiple zeros, it should be noticed that the expression \hyperref[eq:2.11]{(11)} does not appear in reduced form.

Greater unity is achieved if we let the variable z as well as the values $R(z)$ range over the extended plane. We may define $R(\infty)$ as the limit of $R(z)$ as $z \to \infty$, but this definition would not determine the order of a zero or pole at $\infty$. It is therefore preferable to consider the function $R(1/z)$, which we can rewrite as a rational function $R_{1}(z)$, and set
\begin{equation*}
R (\infty) = R _ {1} (0).
\end{equation*}
If $R_{1}(0) = 0$ or $\infty$, the order of the zero or pole at $\infty$ is defined as the order of the zero or pole of $R_{1}(z)$ at the origin.

With the notation
\begin{equation*}
R (z) = \frac {a _ {0} + a _ {1} z + \cdots + a _ {n} z ^ {n}}{b _ {0} + b _ {1} z + \cdots + b _ {m} z ^ {m}}
\end{equation*}
we obtain
\begin{equation*}
R _ {1} (z) = z ^ {m - n} \frac {a _ {0} z ^ {n} + a _ {1} z ^ {n - 1} + \cdots + a _ {n}}{b _ {0} z ^ {m} + b _ {1} z ^ {m - 1} + \cdots + b _ {m}}
\end{equation*}
where the power $z^{m-n}$ belongs either to the numerator or to the denominator. Accordingly, if m > n $R(z)$ has a zero of order m - n at $\infty$, if m < n the point at $\infty$ is a pole of order n - m, and if m = n
\begin{equation*}
R (\infty) = a _ {n} / b _ {m} \neq 0, \infty .
\end{equation*}
% Original PDF physical page 47
\index{Order, algebraic!of a pole|)}\sourcepage{47}
We\index{Fraction, partial|(}\index{Order, algebraic!of rational function|(}\index{Partial fraction|(} can now count the total number of zeros and poles in the extended plane. The count shows that the number of zeros, including those at $\infty$, is equal to the greater of the numbers m and n. The number of poles is the same. This common number of zeros and poles is called the order of the rational function\index{Rational function}.

If a is any constant, the function $R(z) - a$ has the same poles as $R(z)$, and consequently the same order. The zeros of $R(z) - a$ are roots of the equation $R(z) = a$, and if the roots are counted as many times as the order of the zero indicates, we can state the following result:

A rational function $R(z)$ of order $p$ has $p$ zeros and $p$ poles, and every equation $R(z) = a$ has exactly $p$ roots.

A rational function of order 1 is a linear fraction
\begin{equation*}
S (z) = \frac {\alpha z + \beta}{\gamma z + \delta}
\end{equation*}
with $\alpha\delta - \beta\gamma \neq 0$. Such fractions, or linear transformations, will be studied at length in Chap.~\hyperref[chap:3]{3}, Sec.~\hyperref[sec:3.3]{3}. For the moment we note merely that the equation $w = S(z)$ has exactly one root, and we find indeed
\begin{equation*}
z = S ^ {- 1} (w) = \frac {\delta w - \beta}{- \gamma w + \alpha}.
\end{equation*}
The transformations $S$ and $S^{-1}$ are inverse to each other.

The linear transformation $z + a$ is called a parallel translation\index{Parallel translation}, and 1/z is an inversion. The former has a fixed point at $\infty$, the latter interchanges 0 and $\infty$.

Every rational function has a representation by partial fractions. In order to derive this representation we assume first that $R(z)$ has a pole at $\infty$. We carry out the division of $P(z)$ by $Q(z)$ until the degree of the remainder is at most equal to that of the denominator. The result can be written in the form
\begin{equation*}
R (z) = G (z) + H (z)\tag{12}\label{eq:2.12}
\end{equation*}
where $G(z)$ is a polynomial without constant term, and $H(z)$ is finite at $\infty$. The degree of $G(z)$ is the order of the pole at $\infty$, and the polynomial $G(z)$ is called the singular part of $R(z)$ at $\infty$.

Let the distinct finite poles of $R(z)$ be denoted by $\beta_1, \beta_2, \ldots, \beta_q$. The function $R\left(\beta_j + \frac{1}{\zeta}\right)$ is a rational function of $\zeta$ with a pole at $\zeta = \infty$. By use of the decomposition \hyperref[eq:2.12]{(12)} we can write
\begin{equation*}
R \left(\beta_ {j} + \frac {1}{\zeta}\right) = G _ {j} (\zeta) + H _ {j} (\zeta),
\end{equation*}
% Original PDF physical page 48
\index{Fraction, partial|)}\index{Order, algebraic!of rational function|)}\index{Partial fraction|)}\sourcepage{48}
or with a change of variable
\begin{equation*}
R (z) = G _ {j} \left(\frac {1}{z - \beta_ {j}}\right) + H _ {j} \left(\frac {1}{z - \beta_ {j}}\right).
\end{equation*}
Here $G_{j}\left(\frac{1}{z - \beta_{j}}\right)$ is a polynomial in $\frac{1}{z - \beta_{j}}$ without constant term, called the singular part of $R(z)$ at $\beta_{j}$. The function $H_{j}\left(\frac{1}{z - \beta_{j}}\right)$ is finite for $z = \beta_{j}$.

Consider now the expression
\begin{equation*}
R (z) - G (z) - \sum_ {j = 1} ^ {q} G _ {j} \left(\frac {1}{z - \beta_ {j}}\right).\tag{13}\label{eq:2.13}
\end{equation*}
This is a rational function\index{Rational function} which cannot have other poles than $\beta_{1}, \beta_{2}, \ldots, \beta_{q}$ and $\infty$. At $z = \beta_{j}$ we find that the two terms which become infinite have a difference $H_{j}\left(\frac{1}{z - \beta_{j}}\right)$ with a finite limit, and the same is true at $\infty$. Therefore \hyperref[eq:2.13]{(13)} has neither any finite poles nor a pole at $\infty$. A rational function without poles must reduce to a constant, and if this constant is absorbed in $G(z)$ we obtain
\begin{equation*}
R (z) = G (z) + \sum_ {j = 1} ^ {q} G _ {j} \left(\frac {1}{z - \beta_ {j}}\right).\tag{14}\label{eq:2.14}
\end{equation*}
This representation is well known from the calculus where it is used as a technical device in integration theory. However, it is only with the introduction of complex numbers that it becomes completely successful.

\subsection*{EXERCISES}
\originalnumber{1}{exercise:2.1.4.2.1}
1. Use the method of the text to develop
\begin{equation*}
\frac {z ^ {4}}{z ^ {3} - 1} \quad \text { and } \quad \frac {1}{z (z + 1) ^ {2} (z + 2) ^ {3}}
\end{equation*}
in partial fractions.

\originalnumber{2}{exercise:2.1.4.2.2}
2. If $Q$ is a polynomial with distinct roots $\alpha_{1}, \ldots, \alpha_{n}$, and if $P$ is a polynomial of degree $< n$, show that
\begin{equation*}
\frac {P (z)}{Q (z)} = \sum_ {k = 1} ^ {n} \frac {P (\alpha_ {k})}{Q ^ {\prime} (\alpha_ {k}) (z - \alpha_ {k})}.
\end{equation*}
\originalnumber{3}{exercise:2.1.4.2.3}
3. Use the formula in the preceding exercise to prove that there exists a unique polynomial $P$ of degree $< n$ with given values $c_k$ at the points $\alpha_k$ (Lagrange's interpolation polynomial).


% Original PDF physical page 49
\sourcepage{49}
\originalnumber{4}{exercise:2.1.4.2.4}
4.\index{Convergent sequence|(}\index{Sequence!convergent|(}\index{Sequence!divergent|(}\index{Sequence!fundamental|(} What is the general form of a rational function\index{Rational function} which has absolute value 1 on the circle $|z| = 1$? In particular, how are the zeros and poles related to each other?

\originalnumber{5}{exercise:2.1.4.2.5}
5. If a rational function is real on $|z| = 1$, how are the zeros and poles situated?

\originalnumber{6}{exercise:2.1.4.2.6}
6. If $R(z)$ is a rational function of order $n$, how large and how small can the order of $R'(z)$ be?

\setcounter{section}{1}\section{Elementary Theory Of Power Series}\label{sec:2.2}
Polynomials and rational functions are very special analytic functions. The easiest way to achieve greater variety is to form limits. For instance, the sum of a convergent series is such a limit. If the terms are functions of a variable, so is the sum, and if the terms are analytic functions, chances are good that the sum will also be analytic.

Of all series with analytic terms the power series with complex coefficients are the simplest. In this section we study only the most elementary properties of power series. A strong motivation for taking up this study when we are not yet equipped to prove the most general properties (those that depend on integration) is that we need power series to construct the exponential function (Sec. 3).

\setcounter{section}{2}\setcounter{subsection}{0}\subsection{Sequences}\label{sec:2.2.1}
The sequence $\{a_{n}\}_{1}^{\infty}$ has the limit A if to every $\varepsilon > 0$ there exists an $n_{0}$ such that $|a_{n} - A| < \varepsilon$ for $n \geq n_{0}$. A sequence with a finite limit is said to be convergent, and any sequence which does not converge is divergent. If $\lim_{n \to \infty} a_{n} = \infty$, the sequence may be said to diverge to infinity.

Only in rare cases can the convergence be proved by exhibiting the limit, so it is extremely important to make use of a method that permits proof of the existence of a limit even when it cannot be determined explicitly. The test that serves this purpose bears the name of Cauchy. A sequence will be called fundamental, or a Cauchy sequence\index{Cauchy sequence}\index{Sequence!Cauchy}, if it satisfies the following condition: given any $\varepsilon > 0$ there exists an $n_{0}$ such that $|a_{n} - a_{m}| < \varepsilon$ whenever $n \geq n_{0}$ and $m \geq n_{0}$. The test reads:

A sequence is convergent if and only if it is a Cauchy sequence.

The necessity is immediate. If $a_{n} \rightarrow A$ we can find $n_{0}$ such that $|a_{n} - A| < \varepsilon/2$ for $n \geq n_{0}$. For $m, n \geq n_{0}$ it follows by the triangle inequality that $|a_{n} - a_{m}| \leq |a_{n} - A| + |a_{m} - A| < \varepsilon$.

The sufficiency is closely connected with the definition of real numbers, and one way in which real numbers can be introduced is indeed to postulate the sufficiency of Cauchy's condition. However, we wish to use only the property that every bounded monotone sequence of real numbers has a limit.


% Original PDF physical page 50
\index{Convergent sequence|)}\index{Rational function|)}\index{Sequence!convergent|)}\index{Sequence!divergent|)}\index{Sequence!fundamental|)}\sourcepage{50}
The\index{Fundamental sequence|(}\index{Least upper bound (l.u.b.)|(}\index{Limes inferior (symbol, symbol)@Limes inferior ($\liminf$, $\underline{\lim}$)|(}\index{Limes superior (symbol, symbol)@Limes superior ($\limsup$, $\overline{\lim}$)|(} real and imaginary parts of a Cauchy sequence are again Cauchy sequences, and if they converge, so does the original sequence. For this reason we need to prove the sufficiency only for real sequences. We use the opportunity to recall the notions of limes superior and limes inferior. Given a real sequence $\{\alpha_n\}_1^\infty$ we shall set $a_n = \max \{\alpha_1, \ldots, \alpha_n\}$, that is, $a_n$ is the greatest of the numbers $\alpha_1, \ldots, \alpha_n$. The sequence $\{a_n\}_1^\infty$ is nondecreasing; hence it has a limit $A_1$ which is finite or equal to $+\infty$. The number $A_1$ is known as the least upper bound or supremum (l.u.b. or sup) of the numbers $\alpha_n$; indeed, it is the least number which is $\geq$ all $\alpha_n$. Construct in the same way the least upper bound $A_k$ of the sequence $\{\alpha_n\}_k^\infty$ obtained from the original sequence by deleting $\alpha_1, \ldots, \alpha_{k-1}$. It is clear that $\{A_k\}$ is a nonincreasing sequence, and we denote its limit by $A$. It may be finite, $+\infty$, or $-\infty$. In any case we write
\begin{equation*}
A = \lim _ {n \rightarrow \infty} \sup \alpha_ {n}.
\end{equation*}
It is easy to characterize the limes superior by its properties. If A is finite and $\varepsilon > 0$ there exists an $n_{0}$ such that $A_{n_{0}} < A + \varepsilon$, and it follows that $\alpha_{n} \leq A_{n_{0}} < A + \varepsilon$ for $n \geq n_{0}$. In the opposite direction, if $\alpha_{n} \leq A - \varepsilon$ for $n \geq n_{0}$, then $A_{n_{0}} \leq A - \varepsilon$, which is impossible. In other words, there are arbitrarily large n for which $\alpha_{n} > A - \varepsilon$. If $A = +\infty$ there are arbitrarily large $\alpha_{n}$, and $A = -\infty$ if and only if $\alpha_{n}$ tends to $-\infty$. In all cases there cannot be more than one number A with these properties.

The limes inferior can be defined in the same manner with inequalities reversed. It is quite clear that the limes inferior and limes superior will be equal if and only if the sequence converges to a finite limit or diverges to $+\infty$ or to $-\infty$. The notations are frequently simplified to $\overline{\lim}$ and $\underline{\lim}$. The reader should prove the following relations:
\begin{equation*}
\begin{array}{l} \underline {{\lim}} \alpha_ {n} + \underline {{\lim}} \beta_ {n} \leq \underline {{\lim}} (\alpha_ {n} + \beta_ {n}) \leq \underline {{\lim}} \alpha_ {n} + \overline {{\lim}} \beta_ {n} \\ \underline {{\lim}} \alpha_ {n} + \overline {{\lim}} \beta_ {n} \leq \overline {{\lim}} (\alpha_ {n} + \beta_ {n}) \leq \overline {{\lim}} \alpha_ {n} + \overline {{\lim}} \beta_ {n}. \end{array}
\end{equation*}
Now we return to the sufficiency of Cauchy's condition. From $|\alpha_n - \alpha_{n_0}| < \varepsilon$ we obtain $|\alpha_n| < |\alpha_{n_0}| + \varepsilon$ for $n \geq n_0$, and it follows that $A = \varlimsup \alpha_n$ and $a = \underline{\lim}\alpha_n$ are both finite. If $a \neq A$ choose
\begin{equation*}
\varepsilon = \frac {(A - a)}{3}
\end{equation*}
and determine a corresponding $n_0$. By definition of $a$ and $A$ there exists an $\alpha_n < a + \varepsilon$ and an $\alpha_m > A - \varepsilon$ with $m, n \geq n_0$. It follows that $A - a = (A - \alpha_m) + (\alpha_m - \alpha_n) + (\alpha_n - a) < 3\varepsilon$, contrary to the choice of $\varepsilon$. Hence $a = A$, and the sequence converges.


% Original PDF physical page 51
\index{Fundamental sequence|)}\index{Least upper bound (l.u.b.)|)}\index{Limes inferior (symbol, symbol)@Limes inferior ($\liminf$, $\underline{\lim}$)|)}\index{Limes superior (symbol, symbol)@Limes superior ($\limsup$, $\overline{\lim}$)|)}\sourcepage{51}
\setcounter{section}{2}\setcounter{subsection}{1}\subsection{Series}\index{Absolute convergence|(}\index{Convergence!absolute|(}\label{sec:2.2.2}
A very simple application of Cauchy's condition permits us to deduce the convergence of one sequence from that of another. If it is true that $|b_m - b_n| \leq |a_m - a_n|$ for all pairs of subscripts, the sequence $\{b_n\}$ may be termed a contraction\index{Contraction} of the sequence $\{a_n\}$ (this is not a standard term). Under this condition, if $\{a_n\}$ is a Cauchy sequence, so is $\{b_n\}$. Hence convergence of $\{a_n\}$ implies convergence of $\{b_n\}$.

An infinite series is a formal infinite sum
\begin{equation*}
a _ {1} + a _ {2} + \dots + a _ {n} + \dots .\tag{15}\label{eq:2.15}
\end{equation*}
Associated with this series is the sequence of its partial sums
\begin{equation*}
s _ {n} = a _ {1} + a _ {2} + \dots + a _ {n}.
\end{equation*}
The series is said to converge if and only if the corresponding sequence is convergent, and if this is the case the limit of the sequence is the sum of the series.

Applied to a series Cauchy's convergence test yields the following condition: The series \hyperref[eq:2.15]{(15)} converges if and only if to every $\varepsilon > 0$ there exists an $n_0$ such that $|a_n + a_{n+1} + \cdots + a_{n+p}| < \varepsilon$ for all $n \geq n_0$ and $p \geq 0$. For $p = 0$ we find in particular that $|a_n| < \varepsilon$. Hence the general term of a convergent series tends to zero. This condition is necessary, but of course not sufficient.

If a finite number of the terms of the series \hyperref[eq:2.15]{(15)} are omitted, the new series converges or diverges together with \hyperref[eq:2.15]{(15)}. In the case of convergence, let $R_{n}$ be the sum of the series which begins with the term $a_{n+1}$. Then the sum of the whole series is $S = s_{n} + R_{n}$.

The series \hyperref[eq:2.15]{(15)} can be compared with the series
\begin{equation*}
\left| a _ {1} \right| + \left| a _ {2} \right| + \dots + \left| a _ {n} \right| + \dots\tag{16}\label{eq:2.16}
\end{equation*}
formed by the absolute values of the terms. The sequence of partial sums of \hyperref[eq:2.15]{(15)} is a contraction of the sequence corresponding to \hyperref[eq:2.16]{(16)}, for $|a_{n} + a_{n+1} + \cdots + a_{n+p}| \leq |a_{n}| + |a_{n+1}| + \cdots + |a_{n+p}|$. Therefore, convergence of \hyperref[eq:2.16]{(16)} implies that the original series \hyperref[eq:2.15]{(15)} is convergent. A series with the property that the series formed by the absolute values of the terms converges is said to be absolutely convergent.

\setcounter{section}{2}\setcounter{subsection}{2}\subsection{Uniform Convergence}\index{Convergence!uniform|(}\index{Uniform convergence|(}\index{Convergence!uniform}\index{Uniform convergence}\label{sec:2.2.3}
Consider a sequence of functions $f_{n}(x)$, all defined on the same set E. If the sequence of values $\{f_{n}(x)\}$ converges for every x that belongs to E, then the limit $f(x)$ is again a function on E. By definition, if $\varepsilon > 0$ and x belongs to E there exists an $n_{0}$ such that $|f_{n}(x) - f(x)| < \varepsilon$ for $n \geq n_{0}$, but $n_{0}$ is allowed to depend on x.


% Original PDF physical page 52
\index{Absolute convergence|)}\index{Convergence!absolute|)}\sourcepage{52}
For instance, it is true that
\begin{equation*}
\lim _ {n \rightarrow \infty} \left(1 + \frac {1}{n}\right) x = x
\end{equation*}
for all $x$, but in order to have $|(1 + 1/n)x - x| = |x|/n < \varepsilon$ for $n \geq n_0$ it is necessary that $n_0 > |x|/\varepsilon$. Such an $n_0$ exists for every fixed $x$, but the requirement cannot be met simultaneously for all $x$.

We say in this situation that the sequence converges pointwise, but not uniformly. In positive formulation: The sequence $\{f_{n}(x)\}$ converges uniformly to $f(x)$ on the set E if to every $\varepsilon > 0$ there exists an $n_{0}$ such that $|f_{n}(x) - f(x)| < \varepsilon$ for all $n \geq n_{0}$ and all x in E.

The most important consequence of uniform convergence\index{Convergence!uniform}\index{Uniform convergence} is the following:

The limit function of a uniformly convergent sequence of continuous functions is itself continuous.

Suppose that the functions $f_{n}(x)$ are continuous and tend uniformly to $f(x)$ on the set $E$. For any $\varepsilon > 0$ we are able to find an $n$ such that $|f_{n}(x) - f(x)| < \varepsilon / 3$ for all $x$ in $E$. Let $x_0$ be a point in $E$. Because $f_{n}(x)$ is continuous at $x_0$ we can find $\delta > 0$ such that $|f_{n}(x) - f_{n}(x_0)| < \epsilon / 3$ for all $x$ in $E$ with $|x - x_0| < \delta$. Under the same condition on $x$ it follows that
\begin{equation*}
\begin{aligned}
|f(x)-f(x_0)|&\leq |f(x)-f_n(x)|+|f_n(x)-f_n(x_0)|\\
&\qquad+|f_n(x_0)-f(x_0)|<\varepsilon,
\end{aligned}
\end{equation*}
and we have proved that $f(x)$ is continuous at $x_{0}$.

In the theory of analytic functions we shall find uniform convergence much more important than pointwise convergence. However, in most cases it will be found that the convergence is uniform only on a part of the set on which the functions are originally defined.

Cauchy's necessary and sufficient condition has a counterpart for uniform convergence. We assert:

The sequence $\{f_n(x)\}$ converges uniformly on $E$ if and only if to every $\varepsilon > 0$ there exists an $n_0$ such that $|f_m(x) - f_n(x)| < \varepsilon$ for all $m, n \geq n_0$ and all $x$ in $E$.

The necessity is again trivial. For the sufficiency we remark that the limit function $f(x)$ exists by the ordinary form of Cauchy's test. In the inequality $|f_{m}(x) - f_{n}(x)| < \varepsilon$ we can keep n fixed and let m tend to $\infty$. It follows that $|f(x) - f_{n}(x)| \leq \varepsilon$ for $n \geq n_{0}$ and all x in E. Hence the convergence is uniform.

For practical use the following test is the most applicable: If a sequence of functions $\{f_{n}(x)\}$ is a contraction of a convergent sequence of constants $\{a_{n}\}$, then the sequence $\{f_{n}(x)\}$ is uniformly convergent. The hypothesis means that $|f_{m}(x)-f_{n}(x)|\leq|a_{m}-a_{n}|$ on E, and the con%
% Original PDF physical page 53
\sourcepage{53}%
clusion follows immediately by Cauchy's condition.

In the case of series this criterion, in a somewhat weaker form, becomes particularly simple. We say that a series with variable terms
\begin{equation*}
f _ {1} (x) + f _ {2} (x) + \dots + f _ {n} (x) + \dots
\end{equation*}
has the series with positive terms
\begin{equation*}
a _ {1} + a _ {2} + \dots + a _ {n} + \dots
\end{equation*}
for a majorant\index{Majorant} if it is true that $|f_{n}(x)| \leq Ma_{n}$ for some constant M and for all sufficiently large n; conversely, the first series is a minorant\index{Minorant} of the second. In these circumstances we have
\begin{equation*}
\left| f _ {n} (x) + f _ {n + 1} (x) + \dots + f _ {n + p} (x) \right| \leq M \left(a _ {n} + a _ {n + 1} + \dots + a _ {n + p}\right).
\end{equation*}
Therefore, if the majorant converges, the minorant converges uniformly. This condition is frequently referred to as the Weierstrass M test\index{M test}\index{Weierstrass M test}. It has the slight weakness that it applies only to series which are also absolutely convergent. The general principle of contraction is more complicated, but has a wider range of applicability.

\subsection*{EXERCISES}
\originalnumber{1}{exercise:2.2.3.3.1}
1. Prove that a convergent sequence is bounded.

\originalnumber{2}{exercise:2.2.3.3.2}
2. If $\lim_{n\to \infty}z_n = A$, prove that
\begin{equation*}
\lim _ {n \rightarrow \infty} \frac {1}{n} \left(z _ {1} + z _ {2} + \dots + z _ {n}\right) = A.
\end{equation*}
\originalnumber{3}{exercise:2.2.3.3.3}
3. Show that the sum of an absolutely convergent series does not change if the terms are rearranged.

\originalnumber{4}{exercise:2.2.3.3.4}
4. Discuss completely the convergence and uniform convergence\index{Convergence!uniform}\index{Uniform convergence} of the sequence $\{nz^{n}\}_{1}^{\infty}$.

\originalnumber{5}{exercise:2.2.3.3.5}
5. Discuss the uniform convergence of the series
\begin{equation*}
\sum_ {n = 1} ^ {\infty} \frac {x}{n (1 + n x ^ {2})}
\end{equation*}
for real values of x.

\originalnumber{6}{exercise:2.2.3.3.6}
6. If $U = u_{1} + u_{2} + \cdots$, $V = v_{1} + v_{2} + \cdots$ are convergent series, prove that $UV = u_{1}v_{1} + (u_{1}v_{2} + u_{2}v_{1}) + (u_{1}v_{3} + u_{2}v_{2} + u_{3}v_{1}) + \cdots$ provided that at least one of the series is absolutely convergent. (It is easy if both series are absolutely convergent. Try to arrange the proof so economically that the absolute convergence of the second series is not needed.)


% Original PDF physical page 54
\index{Convergence!uniform|)}\index{Uniform convergence|)}\sourcepage{54}
\setcounter{section}{2}\setcounter{subsection}{3}\subsection{Power Series}\index{Abel's power series theorem|(}\label{sec:2.2.4}
A power series\index{Power series|(}\index{Power series} is of the form
\begin{equation*}
a _ {0} + a _ {1} z + a _ {2} z ^ {2} + \dots + a _ {n} z ^ {n} + \dots\tag{17}\label{eq:2.17}
\end{equation*}
where the coefficients $a_{n}$ and the variable $z$ are complex. A little more generally we may consider series
\begin{equation*}
\sum_ {n = 0} ^ {\infty} a _ {n} (z - z _ {0}) ^ {n}
\end{equation*}
which are power series with respect to the center $z_{0}$, but the difference is so slight that we need not do so in a formal manner.

As an almost trivial example we consider the geometric series\index{Geometric series}
\begin{equation*}
1 + z + z ^ {2} + \dots + z ^ {n} + \dots
\end{equation*}
whose partial sums can be written in the form
\begin{equation*}
1 + z + \dots + z ^ {n - 1} = \frac {1 - z ^ {n}}{1 - z}.
\end{equation*}
Since $z^n \to 0$ for $|z| < 1$ and $|z^n| \geq 1$ for $|z| \geq 1$ we conclude that the geometric series converges to $1 / (1 - z)$ for $|z| < 1$, diverges for $|z| \geq 1$.

It turns out that the behavior of the geometric series is typical. Indeed, we shall find that every power series converges inside a circle and diverges outside the same circle, except that it may happen that the series converges only for z = 0, or that it converges for all values of z. More precisely, we shall prove the following theorem due to Abel\index{Abel N H@Abel, N. H.}:

\begin{originalstatement}{Theorem}{2}\label{thm:2.2}
For every power series \hyperref[eq:2.17]{(17)} there exists a number $R$, $0 \leq R \leq \infty$, called the radius of convergence, with the following properties:
(i) The series converges absolutely for every $z$ with $|z| < R$. If $0 \leq \rho < R$ the convergence is uniform for $|z| \leq \rho$.

(ii) If $|z| > R$ the terms of the series are unbounded, and the series is consequently divergent.

(iii) In $|z| < R$ the sum of the series is an analytic function. The derivative can be obtained by termwise differentiation, and the derived series has the same radius of convergence.
\end{originalstatement}

The circle $|z| = R$ is called the circle of convergence\index{Circle of convergence}\index{Convergence!circle of}; nothing is claimed about the convergence on the circle. We shall show that the assertions in the theorem are true if R is chosen according to the formula
\begin{equation*}
1 / R = \lim _ {n \rightarrow \infty} \sup \sqrt [ n ]{| a _ {n} |}.\tag{18}\label{eq:2.18}
\end{equation*}
% Original PDF physical page 55
\sourcepage{55}
This is known as Hadamard's formula\index{Hadamard's formula} for the radius of convergence.

If $|z| < R$ we can find $\rho$ so that $|z| < \rho < R$. Then $1/\rho > 1/R$, and by the definition of limes superior there exists an $n_0$ such that $|a_n|^{1/n} < 1/\rho$, $|a_n| < 1/\rho^n$ for $n \geq n_0$. It follows that $|a_n z^n| < (|z|/\rho)^n$ for large $n$, so that the power series\index{Power series} \hyperref[eq:2.17]{(17)} has a convergent geometric series as a majorant, and is consequently convergent. To prove the uniform convergence for $|z| \leq \rho < R$ we choose a $\rho'$ with $\rho < \rho' < R$ and find $|a_n z^n| \leq (\rho / \rho')^n$ for $n \geq n_0$. Since the majorant is convergent and has constant terms we conclude by Weierstrass's $M$ test that the power series is uniformly convergent.

If $|z| > R$ we choose $\rho$ so that $R < \rho < |z|$. Since $1 / \rho < 1 / R$ there are arbitrarily large $n$ such that $|a_n|^{1 / n} > 1 / \rho$, $|a_n| > 1 / \rho^n$. Thus $|a_n z^n| > (|z| / \rho)^n$ for infinitely many $n$, and the terms are unbounded.

The derived series $\sum_{1}^{\infty} na_n z^{n-1}$ has the same radius of convergence, because $\sqrt[n]{n} \to 1$. Proof: Set $\sqrt[n]{n} = 1 + \delta_n$. Then $\delta_n > 0$, and by use of the binomial theorem $n = (1 + \delta_n)^n > 1 + \frac{1}{2} n(n - 1)\delta_n^2$. This gives $\delta_n^2 < 2/n$, and hence $\delta_n \to 0$.

For $|z| < R$ we shall write
\begin{equation*}
f (z) = \sum_ {0} ^ {\infty} a _ {n} z ^ {n} = s _ {n} (z) + R _ {n} (z)
\end{equation*}
where
\begin{equation*}
s _ {n} (z) = a _ {0} + a _ {1} z + \dots + a _ {n - 1} z ^ {n - 1}, R _ {n} (z) = \sum_ {k = n} ^ {\infty} a _ {k} z ^ {k},
\end{equation*}
and also
\begin{equation*}
f _ {1} (z) = \sum_ {1} ^ {\infty} n a _ {n} z ^ {n - 1} = \lim _ {n \rightarrow \infty} s _ {n} ^ {\prime} (z).
\end{equation*}
We have to show that $f'(z) = f_1(z)$.

Consider the identity
\begin{equation*}
\begin{aligned}
\frac{f(z)-f(z_0)}{z-z_0}-f_1(z_0)
&=\left(\frac{s_n(z)-s_n(z_0)}{z-z_0}-s_n'(z_0)\right)\\
&\quad+(s_n'(z_0)-f_1(z_0))\\
&\quad+\frac{R_n(z)-R_n(z_0)}{z-z_0},
\end{aligned}\tag{19}\label{eq:2.19}
\end{equation*}
where we assume that $z \neq z_0$ and $|z|, |z_0| < \rho < R$. The last term can be rewritten as
\begin{equation*}
\sum_ {k = n} ^ {\infty} a _ {k} (z ^ {k - 1} + z ^ {k - 2} z _ {0} + \dots + z z _ {0} ^ {k - 2} + z _ {0} ^ {k - 1}),
\end{equation*}
% Original PDF physical page 56
\sourcepage{56}
and we conclude that
\begin{equation*}
\left| \frac {R _ {n} (z) - R _ {n} (z _ {0})}{z - z _ {0}} \right| \leq \sum_ {k = n} ^ {\infty} k | a _ {k} | \rho^ {k - 1}.
\end{equation*}
The expression on the right is the remainder term in a convergent series. Hence we can find $n_{0}$ such that
\begin{equation*}
\left| \frac {R _ {n} (z) - R _ {n} (z _ {0})}{z - z _ {0}} \right| <   \frac {\varepsilon}{3}
\end{equation*}
for $n \geq n_0$.

There is also an $n_1$ such that $|s_n'(z_0) - f_1(z_0)| < \varepsilon / 3$ for $n \geq n_1$. Choose a fixed $n \geq n_0, n_1$. By the definition of derivative we can find $\delta > 0$ such that $0 < |z - z_0| < \delta$ implies
\begin{equation*}
\left| \frac {s _ {n} (z) - s _ {n} (z _ {0})}{z - z _ {0}} - s _ {n} ^ {\prime} (z _ {0}) \right| <   \frac {\varepsilon}{3}.
\end{equation*}
When all these inequalities are combined it follows by \hyperref[eq:2.19]{(19)} that
\begin{equation*}
\left| \frac {f (z) - f (z _ {0})}{z - z _ {0}} - f _ {1} (z _ {0}) \right| <   \varepsilon
\end{equation*}
when $0 < |z - z_0| < \delta$. We have proved that $f'(z_0)$ exists and equals $f_1(z_0)$.

Since the reasoning can be repeated we have in reality proved much more: A power series\index{Power series} with positive radius of convergence has derivatives of all orders, and they are given explicitly by
\begin{equation*}
\begin{array}{l} f (z) = a _ {0} + a _ {1} z + a _ {2} z ^ {2} + \cdot \cdot \cdot \\ f ^ {\prime} (z) = a _ {1} + 2 a _ {2} z + 3 a _ {3} z ^ {2} + \cdot \cdot \cdot \\ f ^ {\prime \prime} (z) = 2 a _ {2} + 6 a _ {3} z + 12 a _ {4} z ^ {2} + \cdot \cdot \cdot \\ \cdot \cdot \cdot . \end{array}
\end{equation*}
\begin{equation*}
f ^ {(k)} (z) = k! a _ {k} + \frac {(k + 1) !}{1 !} a _ {k + 1} z + \frac {(k + 2) !}{2 !} a _ {k + 2} z ^ {2} + \dots
\end{equation*}
In particular, if we look at the last line we see that $a_{k} = f^{(k)}(0) / k!$, and the power series becomes
\begin{equation*}
f (z) = f (0) + f ^ {\prime} (0) z + \frac {f ^ {\prime \prime} (0)}{2 !} z ^ {2} + \dots + \frac {f ^ {(n)} (0)}{n !} z ^ {n} + \dots
\end{equation*}
This is the familiar Taylor-Maclaurin development, but we have proved it only under the assumption that $f(z)$ has a power series development. We do know that the development is uniquely determined, if it exists, but the main part is still missing, namely that every analytic function has a Taylor development.


% Original PDF physical page 57
\sourcepage{57}
\subsection*{EXERCISES}
\originalnumber{1}{exercise:2.2.4.4.1}
1. Expand $(1 - z)^{-m}$, m a positive integer, in powers of z.

\originalnumber{2}{exercise:2.2.4.4.2}
2. Expand $\frac{2z + 3}{z + 1}$ in powers of $z - 1$. What is the radius of convergence?

\originalnumber{3}{exercise:2.2.4.4.3}
3. Find the radius of convergence of the following power series\index{Power series}:
\begin{equation*}
\sum n ^ {p} z ^ {n}, \sum \frac {z ^ {n}}{n !}, \sum n! z ^ {n}, \sum q ^ {n ^ {2}} z ^ {n} (| q | <   1), \sum z ^ {n!}
\end{equation*}
\originalnumber{4}{exercise:2.2.4.4.4}
4. If $\Sigma a_{n}z^{n}$ has radius of convergence $R$, what is the radius of convergence of $\Sigma a_{n}z^{2n}$? of $\Sigma a_{n}^{2}z^{n}$?

\originalnumber{5}{exercise:2.2.4.4.5}
5. If $f(z) = \Sigma a_{n}z^{n}$, what is $\Sigma n^{3}a_{n}z^{n}$?

\originalnumber{6}{exercise:2.2.4.4.6}
6. If $\Sigma a_{n}z^{n}$ and $\Sigma b_{n}z^{n}$ have radii of convergence $R_{1}$ and $R_{2}$, show that the radius of convergence of $\Sigma a_{n}b_{n}z^{n}$ is at least $R_{1}R_{2}$.

\originalnumber{7}{exercise:2.2.4.4.7}
7. If $\lim_{n\to \infty}|a_n| / |a_{n + 1}| = R$, prove that $\Sigma a_n z^n$ has radius of convergence $R$.

\originalnumber{8}{exercise:2.2.4.4.8}
8. For what values of z is
\begin{equation*}
\sum_ {0} ^ {\infty} \left(\frac {z}{1 + z}\right) ^ {n}
\end{equation*}
convergent?

\originalnumber{9}{exercise:2.2.4.4.9}
9. Same question for
\begin{equation*}
\sum_ {0} ^ {\infty} \frac {z ^ {n}}{1 + z ^ {2 n}}.
\end{equation*}
\setcounter{section}{2}\setcounter{subsection}{4}\subsection{Abel's Limit Theorem}\index{Abel's limit theorem|(}\index{Abel's limit theorem}\label{sec:2.2.5}
There is a second theorem of Abel's which refers to the case where a power series converges at a point of the circle of convergence. We lose no generality by assuming that R = 1 and that the convergence takes place at z = 1.

\begin{originalstatement}{Theorem}{3}\label{thm:2.3}
If $\sum_{0}^{\infty} a_n$ converges, then $f(z) = \sum_{0}^{\infty} a_n z^n$ tends to $f(1)$ as $z$ approaches 1 in such a way that $|1 - z| / (1 - |z|)$ remains bounded.
\end{originalstatement}
Remark. Geometrically, the condition means that z stays in an angle $<180^{\circ}$ with vertex 1, symmetrically to the part $(-\infty,1)$ of the real axis. It is customary to say that the approach takes place in a Stolz angle\index{Stolz angle}.

Proof. We may assume that $\sum_{0}^{\infty} a_n = 0$, for this can be attained by adding
% Original PDF physical page 58
\index{Abel's power series theorem|)}\sourcepage{58}
a constant to $a_{0}$. We write $s_{n} = a_{0} + a_{1} + \cdots + a_{n}$ and make use of the identity (summation by parts)
\begin{equation*}
\begin{aligned}
s_n(z)&=a_0+a_1z+\dots+a_nz^n\\
&=s_0+(s_1-s_0)z+\dots+(s_n-s_{n-1})z^n\\
&=s_0(1-z)+s_1(z-z^2)+\dots\\
&\qquad+s_{n-1}(z^{n-1}-z^n)+s_nz^n\\
&=(1-z)(s_0+s_1z+\dots+s_{n-1}z^{n-1})+s_nz^n.
\end{aligned}
\end{equation*}
But $s_n z^n \to 0$, so we obtain the representation
\begin{equation*}
f (z) = (1 - z) \sum_ {0} ^ {\infty} s _ {n} z ^ {n}.
\end{equation*}
We are assuming that $|1 - z| \leq K(1 - |z|)$, say, and that $s_n \to 0$. Choose $m$ so large that $|s_n| < \varepsilon$ for $n \geq m$. The remainder of the series $\Sigma s_n z^n$, from $n = m$ on, is then dominated by the geometric series $\varepsilon \sum_{m}^{\infty} |z|^n = \varepsilon |z|^m / (1 - |z|) < \varepsilon / (1 - |z|)$. It follows that
\begin{equation*}
| f (z) | \leq | 1 - z | \left| \sum_ {0} ^ {m - 1} s _ {k} z ^ {k} \right| + K \varepsilon .
\end{equation*}
The first term on the right can be made arbitrarily small by choosing z sufficiently close to 1, and we conclude that $f(z) \to 0$ when $z \to 1$ subject to the stated restriction.

\setcounter{section}{2}\section{The Exponential And Trigonometric Functions}\label{sec:2.3}
The person who approaches calculus exclusively from the point of view of real numbers will not expect any relationship between the exponential function\index{Exponential function|(}\index{Function!exponential|(}\index{Exponential function}\index{Function!exponential} $e^{x}$ and the trigonometric functions $\cos x$ and $\sin x$. Indeed, these functions seem to be derived from completely different sources and with different purposes in mind. He will notice, no doubt, a similarity between the Taylor developments of these functions, and if willing to use imaginary arguments he will be able to derive Euler\index{Euler L@Euler, L.}'s formula $e^{ix} = \cos x + i \sin x$ as a formal identity. But it took the genius of a Gauss to analyze its full depth.

With the preparation given in the preceding section it will be easy to define $e^{z}$, $\cos z$ and $\sin z$ for complex z, and to derive the relations between these functions. At the same time we can define the logarithm as the inverse function of the exponential, and the logarithm leads in turn to the correct definition of the argument of a complex number, and hence to the nongeometric definition of angle.

\setcounter{section}{3}\setcounter{subsection}{0}\subsection{The Exponential}\label{sec:2.3.1}
We may begin by defining the exponential function as the solution of the differential equation


% Original PDF physical page 59
\index{Abel's limit theorem|)}\index{Power series|)}\sourcepage{59}
\begin{equation*}
f ^ {\prime} (z) = f (z)\tag{20}\label{eq:2.20}
\end{equation*}
with the initial value $f(0) = 1$. We solve it by setting
\begin{equation*}
\begin{array}{l} f (z) = a _ {0} + a _ {1} z + \dots + a _ {n} z ^ {n} + \dots \\ f ^ {\prime} (z) = a _ {1} + 2 a _ {2} z + \dots + n a _ {n} z ^ {n - 1} + \dots \end{array}
\end{equation*}
If \hyperref[eq:2.20]{(20)} is to be satisfied, we must have $a_{n-1} = na_n$, and the initial condition gives $a_0 = 1$. It follows by induction that $a_n = 1/n!$.

The solution is denoted by $e^{z}$ or $\exp z$, depending on purely typographical considerations. We must show of course that the series
\begin{equation*}
e ^ {z} = 1 + \frac {z}{1 !} + \frac {z ^ {2}}{2 !} + \dots + \frac {z ^ {n}}{n !} + \dots\tag{21}\label{eq:2.21}
\end{equation*}
converges. It does so in the whole plane, for $\sqrt[n]{n!} \to \infty$ (proof by the reader).

It is a consequence of the differential equation that $e^{z}$ satisfies the addition theorem\index{Addition theorem}
\begin{equation*}
e ^ {a + b} = e ^ {a} \cdot e ^ {b}.\tag{22}\label{eq:2.22}
\end{equation*}
Indeed, we find that $D(e^{z} \cdot e^{c-z}) = e^{z} \cdot e^{c-z} + e^{z} \cdot (-e^{c-z}) = 0$. Hence $e^{z} \cdot e^{c-z}$ is a constant. The value of the constant is found by setting z = 0. We conclude that $e^{z} \cdot e^{c-z} = e^{c}$, and \hyperref[eq:2.22]{(22)} follows for z = a, $c = a + b$.

Remark. We have used the fact that $f(z)$ is constant if $f'(z)$ is identically zero. This is certainly so if f is defined in the whole plane. For if $f = u + iv$ we obtain $\frac{\partial u}{\partial x} = \frac{\partial u}{\partial y} = \frac{\partial v}{\partial x} = \frac{\partial v}{\partial y} = 0$, and the real version of the theorem shows that f is constant on every horizontal and every vertical line.

As a particular case of the addition theorem $e^{z} \cdot e^{-z} = 1$. This shows that $e^{z}$ is never zero. For real x the series development \hyperref[eq:2.21]{(21)} shows that $e^{x} > 1$ for x > 0, and since $e^{x}$ and $e^{-x}$ are reciprocals, $0 < e^{x} < 1$ for x < 0. The fact that the series has real coefficients shows that $\exp \bar{z}$ is the complex conjugate of $\exp z$. Hence $|e^{iy}|^{2} = e^{iy} \cdot e^{-iy} = 1$, and $|e^{x+iy}| = e^{x}$.

\setcounter{section}{3}\setcounter{subsection}{1}\subsection{The Trigonometric Functions}\label{sec:2.3.2}
The trigonometric functions\index{Trigonometric functions|(}\index{Trigonometric functions} are defined by
\begin{equation*}
\cos z = \frac {e ^ {i z} + e ^ {- i z}}{2}, \sin z = \frac {e ^ {i z} - e ^ {- i z}}{2 i}.\tag{23}\label{eq:2.23}
\end{equation*}
Substitution in \hyperref[eq:2.21]{(21)} shows that they have the series developments
\begin{equation*}
\cos z = 1 - \frac {z ^ {2}}{2 !} + \frac {z ^ {4}}{4 !} - \dots
\end{equation*}
% Original PDF physical page 60
\sourcepage{60}
\begin{equation*}
\sin z = z - \frac {z ^ {3}}{3 !} + \frac {z ^ {5}}{5 !} - \dots
\end{equation*}
For real z they reduce to the familiar Taylor developments of $\cos x$ and $\sin x$, with the significant difference that we have now redefined these functions without use of geometry.

From \hyperref[eq:2.23]{(23)} we obtain further Euler\index{Euler L@Euler, L.}'s formula
\begin{equation*}
e ^ {i z} = \cos z + i \sin z
\end{equation*}
as well as the identity
\begin{equation*}
\cos^ {2} z + \sin^ {2} z = 1.
\end{equation*}
It follows likewise that
\begin{equation*}
D \cos z = - \sin z, \quad D \sin z = \cos z.
\end{equation*}
The addition formulas
\begin{equation*}
\begin{array}{r l} & {\cos (a + b) = \cos a \cos b - \sin a \sin b} \\ & {\sin (a + b) = \cos a \sin b + \sin a \cos b} \end{array}
\end{equation*}
are direct consequences of \hyperref[eq:2.23]{(23)} and the addition theorem for the exponential function\index{Exponential function}\index{Function!exponential}.

The other trigonometric functions\index{Trigonometric functions} $\tan z$, $\cot z$, $\sec z$, $\csc z$ are of secondary importance. They are defined in terms of $\cos z$ and $\sin z$ in the customary manner. We find for instance
\begin{equation*}
\tan z = - i \frac {e ^ {i z} - e ^ {- \tau z}}{e ^ {i z} + e ^ {- i z}}.
\end{equation*}
Observe that all the trigonometric functions are rational functions of $e^{iz}$.

\subsection*{EXERCISES}
\originalnumber{1}{exercise:2.3.2.5.1}
1. Find the values of $\sin i$, $\cos i$, $\tan(1+i)$.

\originalnumber{2}{exercise:2.3.2.5.2}
2. The hyperbolic cosine and sine are defined by $\cosh z = (e^{z} + e^{-z})/2$, $\sinh z = (e^{z} - e^{-z})/2$. Express them through $\cos iz$, $\sin iz$. Derive the addition formulas, and formulas for $\cosh 2z$, $\sinh 2z$.

\originalnumber{3}{exercise:2.3.2.5.3}
3. Use the addition formulas to separate $\cos(x + iy)$, $\sin(x + iy)$ in real and imaginary parts.

\originalnumber{4}{exercise:2.3.2.5.4}
4. Show that
\begin{equation*}
| \cos z | ^ {2} = \sinh^ {2} y + \cos^ {2} x = \cosh^ {2} y - \sin^ {2} x = \frac {1}{2} (\cosh 2 y + \cos 2 x)
\end{equation*}
and
\begin{equation*}
| \sin z | ^ {2} = \sinh^ {2} y + \sin^ {2} x = \cosh^ {2} y - \cos^ {2} x = \frac {1}{2} (\cosh 2 y - \cos 2 x).
\end{equation*}
\setcounter{section}{3}\setcounter{subsection}{2}\subsection{The Periodicity}\label{sec:2.3.3}
We say that $f(z)$ has the period $c$ if $f(z + c) = f(z)$


% Original PDF physical page 61
\index{Trigonometric functions|)}\sourcepage{61}
for all z. Thus a period\index{Period|(}\index{Period} of $e^{z}$ satisfies $e^{z+c} = e^{z}$, or $e^{c} = 1$. It follows that $c = i\omega$ with real $\omega$; we prefer to say that $\omega$ is a period of $e^{iz}$. We shall show that there are periods, and that they are all integral multiples of a positive period $\omega_{0}$.

Of the many ways to prove the existence of a period we choose the following: From $D \sin y = \cos y \leq 1$ and $\sin 0 = 0$ we obtain $\sin y < y$ for $y > 0$, either by integration or by use of the mean-value theorem. In the same way $D \cos y = -\sin y > -y$ and $\cos 0 = 1$ gives $\cos y > 1 - y^2 / 2$, which in turn leads to $\sin y > y - y^3 / 6$ and finally to $\cos y < 1 - y^2 / 2 + y^4 / 24$. This inequality shows that $\cos \sqrt{3} < 0$, and therefore there is a $y_0$ between 0 and $\sqrt{3}$ with $\cos y_0 = 0$. Because
\begin{equation*}
\cos^ {2} y _ {0} + \sin^ {2} y _ {0} = 1
\end{equation*}
we have $\sin y_{0} = \pm1$, that is, $e^{iy_{0}} = \pm i$, and hence $e^{4iy_{0}} = 1$. We have shown that $4y_{0}$ is a period.

Actually, it is the smallest positive period. To see this, take $0 < y < y_0$. Then $\sin y > y(1-y^2/6) > y/2 > 0$, which shows that $\cos y$ is strictly decreasing. Because $\sin y$ is positive and $\cos^2 y + \sin^2 y = 1$ it follows that $\sin y$ is strictly increasing, and hence $\sin y < \sin y_0 = 1$. The double inequality $0 < \sin y < 1$ guarantees that $e^{iy}$ is neither $\pm1$ nor $\pm i$. Therefore $e^{4iy} \ne 1$, and $4y_0$ is indeed the smallest positive period. We denote it by $\omega_0$.

Consider now an arbitrary period $\omega$. There exists an integer n such that $n\omega_{0} \leq \omega < (n + 1)\omega_{0}$. If $\omega$ were not equal to $n\omega_{0}$, then $\omega - n\omega_{0}$ would be a positive period $< \omega_{0}$. Since this is not possible, every period must be an integral multiple of $\omega_{0}$.

The smallest positive period of $e^{iz}$ is denoted by $2\pi$.

In the course of the proof we have shown that
\begin{equation*}
e ^ {\pi i / 2} = i, \quad e ^ {\pi i} = - 1, \quad e ^ {2 \pi i} = 1.
\end{equation*}
These equations demonstrate the intimate relationship between the numbers e and $\pi$.

When y increases from 0 to $2\pi$, the point $w = e^{iy}$ describes the unit circle $|w| = 1$ in the positive sense, namely from 1 over i to -1 and back over -i to 1. For every w with $|w| = 1$ there is one and only one y from the half-open interval $0 \leq y < 2\pi$ such that $w = e^{iy}$. All this follows readily from the established fact that cos y is strictly decreasing in the “first quadrant,” that is, between 0 and $\pi/2$.

From an algebraic point of view the mapping $w = e^{iy}$ establishes a homomorphism\index{Homomorphism} between the additive group of real numbers and the multiplicative group of complex numbers with absolute value 1. The kernel\index{Kernel} of the homomorphism is the subgroup formed by all integral multiples $2\pi n$.


% Original PDF physical page 62
\sourcepage{62}
\setcounter{section}{3}\setcounter{subsection}{3}\subsection{The Logarithm}\index{pi@$\pi$|(}\label{sec:2.3.4}
Together with the exponential function\index{Exponential function}\index{Function!exponential} we must also study its inverse function, the logarithm\index{Logarithm|(}\index{Logarithm}. By definition, $z = \log w$ is a root of the equation $e^{z} = w$. First of all, since $e^{z}$ is always $\neq 0$, the number 0 has no logarithm. For $w \neq 0$ the equation $e^{x+iy} = w$ is equivalent to
\begin{equation*}
e ^ {x} = | w |, \quad e ^ {i y} = w / | w |.\tag{24}\label{eq:2.24}
\end{equation*}
The first equation has a unique solution $x = \log |w|$, the real logarithm of the positive number $|w|$. The right-hand member of the second equation \hyperref[eq:2.24]{(24)} is a complex number of absolute value 1. Therefore, as we have just seen, it has one and only one solution in the interval $0 \leq y < 2\pi$. In addition, it is also satisfied by all y that differ from this solution by an integral multiple of $2\pi$. We see that every complex number other than 0 has infinitely many logarithms which differ from each other by multiples of $2\pi i$.

The imaginary part of $\log w$ is also called the argument\index{Argument} of w, $\arg w$, and it is interpreted geometrically as the angle\index{Angle}, measured in radians, between the positive real axis and the half line from 0 through the point w. According to this definition the argument has infinitely many values which differ by multiples of $2\pi$, and
\begin{equation*}
\log w = \log | w | + i \arg w.
\end{equation*}
With a change of notation, if $|z| = r$ and $\arg z = \theta$, then $z = re^{i\theta}$. This notation is so convenient that it is used constantly, even when the exponential function is not otherwise involved.

By convention the logarithm of a positive number shall always mean the real logarithm, unless the contrary is stated. The symbol $a^{b}$, where a and b are arbitrary complex numbers except for the condition $a \neq 0$, is always interpreted as an equivalent of $\exp(b \log a)$. If a is restricted to positive numbers, $\log a$ shall be real, and $a^{b}$ has a single value. Otherwise $\log a$ is the complex logarithm, and $a^{b}$ has in general infinitely many values which differ by factors $e^{2\pi in b}$. There will be a single value if and only if b is an integer n, and then $a^{b}$ can be interpreted as a power of a or $a^{-1}$. If b is a rational number with the reduced form p/q, then $a^{b}$ has exactly q values and can be represented as $\sqrt[q]{a^{p}}$.

The addition theorem of the exponential function clearly implies
\begin{equation*}
\begin{array}{l} \log (z _ {1} z _ {2}) = \log z _ {1} + \log z _ {2} \\ \arg (z _ {1} z _ {2}) = \arg z _ {1} + \arg z _ {2}, \end{array}
\end{equation*}
but only in the sense that both sides represent the same infinite set of complex numbers. If we want to compare a value on the left with a value on the right, then we can merely assert that they differ by a multiple of $2\pi i$ (or $2\pi$ ). (Compare with the remarks in Chap.~\hyperref[chap:1]{1}, Sec.~\hyperref[sec:1.2.1]{2.1}.)


% Original PDF physical page 63
\index{Period|)}\index{pi@$\pi$|)}\sourcepage{63}
Finally we discuss the inverse cosine which is obtained by solving the equation
\begin{equation*}
\cos z = \frac {1}{2} (e ^ {i z} + e ^ {- i z}) = w.
\end{equation*}
This is a quadratic equation in $e^{iz}$ with the roots
\begin{equation*}
e ^ {i z} = w \pm \sqrt {w ^ {2} - 1},
\end{equation*}
and consequently
\begin{equation*}
z = \operatorname{arc} \cos w = - i \log (w \pm \sqrt {w ^ {2} - 1}).
\end{equation*}
We can also write these values in the form
\begin{equation*}
\operatorname{arc} \cos w = \pm i \log (w + \sqrt {w ^ {2} - 1}),
\end{equation*}
for $w + \sqrt{w^{2} - 1}$ and $w - \sqrt{w^{2} - 1}$ are reciprocal numbers. The infinitely many values of arc cos w reflect the evenness and periodicity of cos z. The inverse sine is most easily defined by
\begin{equation*}
\operatorname{arc} \sin w = \frac {\pi}{2} - \operatorname{arc} \cos w.
\end{equation*}
It is worth emphasizing that in the theory of complex analytic functions all elementary transcendental functions can thus be expressed through $e^{z}$ and its inverse $\log z$. In other words, there is essentially only one elementary transcendental function.

\subsection*{EXERCISES}
\originalnumber{1}{exercise:2.3.4.6.1}
1. For real y, show that every remainder in the series for $\cos y$ and $\sin y$ has the same sign as the leading term (this generalizes the inequalities used in the periodicity proof, Sec.~\hyperref[sec:2.3.3]{3.3}).

\originalnumber{2}{exercise:2.3.4.6.2}
2. Prove, for instance, that $3 < \pi < 2\sqrt{3}$.

\originalnumber{3}{exercise:2.3.4.6.3}
3. Find the value of $e^z$ for $z = -\frac{\pi i}{2}, \frac{3}{4} \pi i, \frac{2}{3} \pi i$.

\originalnumber{4}{exercise:2.3.4.6.4}
4. For what values of $z$ is $e^z$ equal to 2, -1, i, -i/2, -1 - i, 1 + 2i?

\originalnumber{5}{exercise:2.3.4.6.5}
5. Find the real and imaginary parts of $\exp (e^z)$.

\originalnumber{6}{exercise:2.3.4.6.6}
6. Determine all values of $2^{i}, i^{i}, (-1)^{2i}$.

\originalnumber{7}{exercise:2.3.4.6.7}
7. Determine the real and imaginary parts of $z^{z}$.

\originalnumber{8}{exercise:2.3.4.6.8}
8. Express arc tan w in terms of the logarithm\index{Logarithm}.

\originalnumber{9}{exercise:2.3.4.6.9}
9. Show how to define the “angles” in a triangle, bearing in mind that they should lie between 0 and $\pi$. With this definition, prove that the sum of the angles is $\pi$.

\originalnumber{10}{exercise:2.3.4.6.10}
10. Show that the roots of the binomial equation $z^{n} = a$ are the vertices of a regular polygon (equal sides and angles).


% Original PDF physical page 64
\index{Complex function|)}\index{Exponential function|)}\index{Function!exponential|)}\sourcepage{64}
\index{Function!complex|)}\index{Logarithm|)}